\chapter{Discussion}\label{ch:discussion}

The results presented in \Cref{ch:evaluation} depict Vinculum as possessing a
shorter running time and scaling at least as good as the other participating programs.

While this is the case for the scenarios measured by the benchmarks it is important
to remember the differences in internal complexity as described in \Cref{sec:usage}.
Since Vinculum only concerns itself with performing the fossil collection and deletion
steps the overall performance of programs using it is strongly influenced by their
individual implementations of the repository and backup creation.

A more conclusive performance evaluation needs to measure the behavior in real-world
situations which includes the creation of backups and remote repositories.
While the benchmarks performed in \Cref{ch:evaluation} attempt to provide a
level playing field the vinculum-benchmark program described in \Cref{sec:repo-implementation}
provides only primitive options for creating backups while not dealing with concerns
such as privacy or reducing possible network access through caching and is therefore
unusable as a real backup solution.

This also highlights the amount of work required to create a backup solution
comparable with the existing ones.
While Vinculum extends the Lock-Free Deduplication algorithm with features
such as client management (\Cref{sec:adding-clients}, \Cref{sec:removing-clients})
and error recovery (\Cref{sec:abandoned-fossils}, \Cref{sec:unreferenced-chunks})
and provides the building blocks the actual implementation is deferred to its users.
Since this requires the consideration of additional concerns the users have to
implement features such as encryption, compression and caching on their own.

The paper describing the Lock-Free Deduplication algorithm does provide
descriptions of how encryption, compression and caching can be implemented, but
the solutions restrict the way other concerns such as the backup retention policy
or chunk creation are
implemented~\cite[sections `Caching' and `Compression and Encryption']{Duplicacy}.
While the development of Vinculum started with the intent of providing a unified
interface on which these features could be implemented as reusable independent
components it was abandoned for reasons explained in \Cref{sec:changes}.

There are also some scenarios which where not discussed in \Cref{ch:design}.
One of them is the repository experiencing catastrophic errors such as a power
failure.
While \Cref{sec:errors} provides solutions to recover from clients experiencing
them and the approach described in \Cref{sec:creation} can be modified to guard
against partially written files being left behind there may be unforeseen interactions
with the fossil collection and deletion steps.
